\chapter{Empirical transitions and concentration events}
\label{chap:empirical}

This chapter defines the statistics that any learner in the episode environment computes from
the observed episodes (Section~2 of \cite{essakine2026tight}, ``Empirical MDP''), the
exploration rates and the three concentration events of Appendix~A of the paper, and proves
Lemma~5 of the paper: the good event has probability at least $1 - \delta$.

\textbf{Standing setting.} $\mathcal M$ is a finite-horizon MDP with a reward function $r$
(Definition~\ref{def:reward_fn}), $s_1$ an initial state, and $(X, Y)$ is an
algorithm-environment sequence for an \emph{arbitrary} algorithm and the episode environment of
$\mathcal M$ from $s_1$ (Definition~\ref{def:episode_env}), on a probability space
$(\Omega, \Prob)$: $X_t = \pi^{t+1}$ is the policy of the episode $t + 1$ and
$Y_t = (s^{t+1}_1, \dots, s^{t+1}_{H+1})$ its state sequence, $a^{t+1}_h := \pi^{t+1}_h(s^{t+1}_h)$.
The \emph{history of the first $t$ episodes} is $\hist_t := (X_i, Y_i)_{i < t}$ and we write
$\mathcal F_t := \sigma(\hist_t, X_t)$ for the $\sigma$-algebra generated by the history of the
first $t$ episodes together with the policy of the episode $t + 1$ ($t \ge 0$; the filtration
$(\mathcal F_t)_{t \ge 0}$ is nondecreasing since $X_t$ is part of $\hist_{t+1}$). Nothing in this
chapter depends on the algorithm; Chapter~\ref{chap:algorithm} specializes to Entropic-BPI,
for which $X_t$ is a function of $\hist_t$ and $\mathcal F_t = \sigma(\hist_t)$.

\textbf{Lean remark.} The history of the first $t$ episodes of the run is
\texttt{histAt X Y t ω : Hist Unit (Policy S A H) (Traj S H) t} (\texttt{Run.lean}), LML's
\texttt{history O X Y t} with the trivial observations; all the quantities below are first
defined as functions of a history \texttt{hist : Hist Unit (Policy S A H) (Traj S H) t}
(\texttt{Empirical.lean}, namespace \texttt{Learning.MDP}) and evaluated on the run through
\texttt{histAt}. The $\sigma$-algebra $\mathcal F_t$ is generated by the pair
\texttt{(histAt X Y t, X t)}; LML's history filtration \texttt{IsAlgEnvSeq.filtration n} is
generated by the rounds $0, \dots, n$, i.e.\ by $\hist_{n+1}$, so that $\mathcal F_t$ sits between
\texttt{filtration (t - 1)} and \texttt{filtration t}. The concentration statements are
formulated with conditional distributions (\texttt{HasCondDistrib}) given
\texttt{(histAt X Y t, X t)}, which avoids naming the filtration.

\section{Counts and empirical transitions}
\label{sec:empirical_counts}

\begin{definition}[Visit counts]\label{def:counts}
\lean{Learning.MDP.EmpiricalModel, Learning.MDP.EmpiricalModel.ofEpisodeAt, Essakine2026Tight.stepModel, Essakine2026Tight.visitCount, Essakine2026Tight.histAt, Essakine2026Tight.visitCountAt}\leanok
\uses{def:episode_env}
For $t \ge 0$, $h \in \{1, \dots, H\}$, $(s, a) \in \Ss \times \As$ and $s' \in \Ss$, the
\emph{visit counts} after $t$ episodes are
\[
  n_h^t(s, a) := \sum_{i = 1}^{t} \indic\{(s^i_h, a^i_h) = (s, a)\},
  \qquad
  n_h^t(s, a, s') := \sum_{i = 1}^{t} \indic\{(s^i_h, a^i_h, s^i_{h+1}) = (s, a, s')\},
\]
where $a^i_h = \pi^i_h(s^i_h)$. They are functions of the history $\hist_t$ of the first $t$
episodes (hence $\mathcal F_t$-measurable, and even $\sigma(\hist_t)$-measurable),
$n_h^t(s, a) = \sum_{s'} n_h^t(s, a, s')$, $n_h^0 = 0$, and $n_h^{t+1}(s, a) - n_h^t(s, a) \in \{0, 1\}$.
\end{definition}

\textbf{Lean remark.} \texttt{EmpiricalModel S A := (S × A → ℕ) × (S × A → ℝ) × (S × A → S → ℕ)}
(visit counts, reward sums, transition counts; an additive commutative monoid, so that the
model of a history is a finite sum), \texttt{EmpiricalModel.ofEpisodeAt π τ r h} the model of
the single transition at step \texttt{h} of the episode with policy \texttt{π}, states
\texttt{τ} and rewards \texttt{r}, \texttt{stepModel hist h := ∑ i, ofEpisodeAt (hist i).action (hist i).feedback 0 h}
(the rewards are not observed, hence set to $0$), \texttt{visitCount hist h s a := (stepModel hist h).count (s, a)};
on the run, \texttt{visitCountAt X Y h s a t ω := visitCount (histAt X Y t ω) h s a}.
Lemma~\ref{lem:count_le_episodes} ($n_h^t(s, a) \le t$, $\sum_{s, a} n_h^t(s, a) = t$) is in
Chapter~\ref{chap:pre_mdp}.

\begin{definition}[Empirical transitions]\label{def:emp_trans}
\lean{Essakine2026Tight.empTrans, Essakine2026Tight.empTransAt}\leanok
\uses{def:counts}
The \emph{empirical transition probabilities} after $t$ episodes are
\[
  \phat_h^t(s' \mid s, a) :=
  \begin{cases}
    n_h^t(s, a, s') / n_h^t(s, a) & \text{if } n_h^t(s, a) > 0, \\
    1/S & \text{otherwise},
  \end{cases}
\]
and $\phat_h^t(s, a) := \phat_h^t(\cdot \mid s, a) \in \Sigma_{\Ss}$ denotes the corresponding
probability vector; as in Definition~\ref{def:episodic_mdp}, $(\phat_h^t f)(s, a) := \sum_{s'} \phat_h^t(s' \mid s, a) f(s')$
and $\Var_{\phat_h^t}(f)(s, a) := \phat_h^t\big((f - \phat_h^t f(s, a))^2\big)(s, a)$.
\end{definition}

\textbf{Lean remark.} \texttt{empTrans hist h s a : S → ℝ} (\texttt{if visitCount hist h s a = 0 then fun \_ ↦ (Fintype.card S : ℝ)⁻¹ else (stepModel hist h).empTrans (s, a)}),
\texttt{empTransAt X Y h s a t ω := empTrans (histAt X Y t ω) h s a}. The value $1/S$ for a
never-visited pair is the paper's convention; it never matters, since every use of $\phat$
for a pair with $n = 0$ is replaced by a trivial bound (deviation~1 of the overview).

\begin{lemma}[The empirical transitions form a probability vector]\label{lem:emp_trans_simplex}
\uses{def:emp_trans}
For all $t, h, (s, a)$: $\phat_h^t(\cdot \mid s, a)$ is a probability vector on $\Ss$
(nonnegative entries summing to $1$). Consequently $f \mapsto (\phat_h^t f)(s, a)$ is linear,
positive and monotone, $\phat_h^t \mathbf 1 = 1$, $\abs{(\phat_h^t f)(s, a)} \le \max_{s'} \abs{f(s')}$,
$\Var_{\phat_h^t}(f)(s, a) \ge 0$, and if $f$ takes values in $[m, M]$ then so does $(\phat_h^t f)(s, a)$.
\end{lemma}

\begin{proof}
\uses{def:emp_trans, def:counts}
If $n_h^t(s, a) > 0$, the entries $n_h^t(s, a, s')/n_h^t(s, a)$ are nonnegative and sum to
$\sum_{s'} n_h^t(s, a, s')/n_h^t(s, a) = 1$ (Definition~\ref{def:counts}); otherwise the entries are
$1/S$, which sum to $1$. The consequences hold for every probability vector $q$: linearity,
positivity and $q \mathbf 1 = 1$ are immediate from $q f = \sum_{s'} q(s') f(s')$ with
$q(s') \ge 0$, $\sum q(s') = 1$; the bound and the preservation of $[m, M]$ follow from the
convexity of the combination; $\Var_q(f) = q(f - qf)^2 \ge 0$ by positivity.
\end{proof}

\textbf{Lean remark.} \texttt{isProbVec\_empTrans}, then the \texttt{vecExp}/\texttt{vecVar}
lemmas of \texttt{Vec.lean} (Lemma~\ref{lem:exp_value_monotone}) apply to
\texttt{empTrans hist h s a}. The KL divergence of Definition~\ref{def:event_kl} views this
vector as a measure through \texttt{weightedMeasure : (S → ℝ) → Measure S}
(\texttt{Essakine2026Tight/Mathlib/}), $\sum_{s'} q(s') \delta_{s'}$.

\section{Pseudo-counts}
\label{sec:empirical_pseudo}

\begin{definition}[Pseudo-counts]\label{def:pseudo_counts}
\lean{Essakine2026Tight.pseudoCount}\leanok
\uses{def:counts, def:occupancy, lem:episode_env_law}
The \emph{pseudo-count} of $(s, a)$ at step $h$ after $t$ episodes is the sum of the conditional
probabilities of the visits given the past:
\[
  \nbar_h^t(s, a) := \sum_{i = 1}^{t} p^{\pi^i}_h(s, a),
  \qquad
  p^\pi_h(s, a) := \Prob^\pi_{(s_1, 1)}(S_h = s, A_h = a)
\]
(the occupancy measure of $\pi$, Definition~\ref{def:occupancy}). It is a random variable, a
function of the policies $\pi^1, \dots, \pi^t$ played, hence $\mathcal F_{t-1}$-measurable for
$t \ge 1$ (predictable), and $\nbar_h^0 = 0$.
\end{definition}

\begin{remark}[Deviation 10: random pseudo-counts]\label{rem:empirical_pseudo_counts}
The paper defines the pseudo-counts as the expectations $\E[n_h^t(s, a)]$ (following
\cite{menard2021fast}), and then uses two properties: Lemma~20 of the paper (the Bernoulli
maximal inequality, Theorem~\ref{thm:bernoulli_concentration}), which controls
$n_h^t(s, a)$ by $\sum_{i \le t} \Prob(\text{visit at episode } i \mid \text{past})$, a
\emph{random} sum; and the counting argument of its Appendix~B.1 (Lemma~\ref{lem:sum_counts_bound}),
which telescopes the increments $\nbar_h^{t+1} - \nbar_h^t = p^{\pi^{t+1}}_h(s, a)$, again the random
conditional visit probabilities. Both are properties of the random predictable sum, not of its
expectation (the two differ as soon as the algorithm is adaptive). The blueprint therefore
\emph{defines} the pseudo-counts as the random sums; Lemma~\ref{lem:pseudo_counts_increments}(iv)
shows that the paper's quantity is their expectation. Since the policies are deterministic,
the paper's ``$\nu_{\hat\pi}(u) > 1/2$'' in the lower bound is ``$\hat\pi$ realizes $u$''.
\end{remark}

\textbf{Lean remark.} \texttt{pseudoCount M s₁ X h s a t ω := ∑ i ∈ Finset.range t, occupancy M (X i ω) s₁ h s a}
(the outline's \texttt{pseudoCount X Y P h s a t}, now depending on the MDP through the
occupancy measures and not on the measure \texttt{P}; the expectation version
$\E[n_h^t(s, a)]$ is not used).

\begin{lemma}[Increments of the pseudo-counts]\label{lem:pseudo_counts_increments}
\uses{def:pseudo_counts, def:counts, lem:episode_env_law, def:occupancy, lem:occupancy_sum_one}
For all $t \ge 0$, $h \le H$ and $(s, a)$:
\begin{enumerate}
  \item[(i)] $\nbar_h^{t+1}(s, a) - \nbar_h^t(s, a) = p^{\pi^{t+1}}_h(s, a) \in [0, 1]$, with
        $p^{\pi}_h(s, a) = \indic\{a = \pi_h(s)\}\, \Prob^\pi_{(s_1, 1)}(S_h = s)$;
  \item[(ii)] $0 \le \nbar_h^t(s, a) \le t$ and $\sum_{s, a} \nbar_h^t(s, a) = t$;
  \item[(iii)] $\Prob\big((s^{t+1}_h, a^{t+1}_h) = (s, a) \mid \mathcal F_t\big) = p^{\pi^{t+1}}_h(s, a)$
        almost surely;
  \item[(iv)] $\E[n_h^t(s, a)] = \E[\nbar_h^t(s, a)]$: the paper's pseudo-count is the expectation
        of $\nbar_h^t(s, a)$.
\end{enumerate}
\end{lemma}

\begin{proof}
\uses{def:pseudo_counts, def:counts, lem:episode_env_law, lem:occupancy_sum_one}
(i) is the definition; the formula for $p^\pi_h(s, a)$ and $\sum_{s, a} p^\pi_h(s, a) = 1$ are
Lemma~\ref{lem:occupancy_sum_one}, which gives $p^\pi_h(s, a) \in [0, 1]$. (ii) follows from
(i) by induction on $t$. (iii) By Lemma~\ref{lem:episode_env_law}, the conditional law of
$Y_t = (s^{t+1}_1, \dots, s^{t+1}_{H+1})$ given $\mathcal F_t = \sigma(\hist_t, X_t)$ is the law of
the state sequence of the policy $X_t = \pi^{t+1}$ from $s_1$; the event
$\{(s^{t+1}_h, a^{t+1}_h) = (s, a)\} = \{s^{t+1}_h = s\} \cap \{\pi^{t+1}_h(s) = a\}$ is the
event $\{S_h = s, A_h = a\}$ of this state sequence (with $A_h = \pi^{t+1}_h(S_h)$), whose
probability under $\Prob^{\pi^{t+1}}_{(s_1, 1)}$ is $p^{\pi^{t+1}}_h(s, a)$. (iv) Taking
expectations in (iii) (tower rule), $\E[\indic\{(s^{i}_h, a^{i}_h) = (s, a)\}] = \E[p^{\pi^{i}}_h(s, a)]$
for every $i \ge 1$; sum over $i = 1, \dots, t$.
\end{proof}

\section{Exploration rates}
\label{sec:empirical_rates}

\begin{definition}[Exploration rates]\label{def:rates}
\lean{Essakine2026Tight.alphaKL, Essakine2026Tight.alphaStar, Essakine2026Tight.alphaCnt}\leanok
\uses{def:episodic_mdp}
For $\delta \in (0, 1)$ and $n \in \N$, with $S = \abs{\Ss}$, $A = \abs{\As}$ (Appendix~A of the
paper, eq.~(5)):
\[
  \alpha(n, \delta) := \log\frac{3 S A H}{\delta} + S \log\big(8 e (n + 1)\big),
  \qquad
  \alpha^*(n, \delta) := \log\frac{3 S A H}{\delta} + \log\big(8 e (n + 1)\big),
  \qquad
  \alpha^{\mathrm{cnt}}(\delta) := \log\frac{3 S A H}{\delta} .
\]
We write $\alpha(n)$, $\alpha^*(n)$ when $\delta$ is fixed.
\end{definition}

\textbf{Lean remark.} \texttt{alphaKL S A H δ n}, \texttt{alphaStar S A H δ n},
\texttt{alphaCnt S A H δ} (\texttt{Essakine2026Tight/EV2026/Rates.lean}), functions of the
cardinalities \texttt{S A : ℕ}, of \texttt{H : ℕ} and of \texttt{δ : ℝ}; they are applied
with \texttt{Fintype.card S}, \texttt{Fintype.card A}.

\begin{lemma}[Properties of the rates]\label{lem:rates_props}
\uses{def:rates}
For $\delta \in (0, 1)$ (more generally $\delta \in (0, 1]$):
\begin{enumerate}
  \item[(i)] $0 \le \alpha^{\mathrm{cnt}}(\delta) \le \alpha^*(n, \delta) \le \alpha(n, \delta)$ for all $n$,
        and $\alpha^{\mathrm{cnt}}(\delta) \ge \log 3 > 1$;
  \item[(ii)] $n \mapsto \alpha(n, \delta)$ and $n \mapsto \alpha^*(n, \delta)$ are nondecreasing;
  \item[(iii)] $n \mapsto \alpha(n, \delta)/n$ and $n \mapsto \alpha^*(n, \delta)/n$ are nonincreasing on
        $n \ge 1$: $\alpha(n, \delta)/n \le \alpha(n', \delta)/n'$ for $n \ge n' \ge 1$, and the same for $\alpha^*$;
  \item[(iv)] $\alpha(n, \delta) \le \log(3SAH/\delta) + S \log(8e(n + 1))$ and
        $\alpha^*(n, \delta) \le \log(3SAH/\delta) + \log(8e(n + 1))$ (equalities, recorded as the
        only form in which the rates enter the final bound).
\end{enumerate}
\end{lemma}

\begin{proof}
\uses{def:rates}
(i) $S, A, H \ge 1$ and $\delta \le 1$ give $3SAH/\delta \ge 3$, so $\alpha^{\mathrm{cnt}} \ge \log 3 > 1 > 0$;
$\log(8e(n + 1)) \ge \log(8e) > 0$ and $S \ge 1$ give the two other inequalities.
(ii) $n \mapsto \log(8e(n + 1))$ is nondecreasing. (iii) Write $\alpha(n)/n = c/n + S \log(n + 1)/n$
with $c := \log(3SAH/\delta) + S \log(8e) \ge 0$; $c/n$ is nonincreasing, and $x \mapsto \log(1 + x)/x$
is nonincreasing on $(0, \infty)$ because $x \mapsto \log(1 + x)$ is concave with value $0$ at
$0$ (for $0 < x' \le x$, concavity gives $\log(1 + x') \ge (x'/x)\log(1 + x) + (1 - x'/x)\log 1$).
The same argument applies to $\alpha^*$ with $S$ replaced by $1$. (iv) is the definition.
\end{proof}

\section{Concentration events}
\label{sec:empirical_events}

The three events below are defined for the run $(X, Y)$ and a confidence level
$\delta \in (0, 1)$; they are events of $\Omega$ (countable intersections of measurable sets).

\begin{definition}[KL concentration event]\label{def:event_kl}
\lean{Essakine2026Tight.eventKL, Essakine2026Tight.klThreshold}\leanok
\uses{def:emp_trans, def:counts, def:rates}
\[
  \evKL := \Big\{ \forall t \in \N,\ \forall h \le H,\ \forall (s, a) \in \Ss \times \As :\
    \KL\big(\phat_h^t(s, a) \,\big\|\, p_h(s, a)\big) \le \frac{\alpha(n_h^t(s, a), \delta)}{n_h^t(s, a)} \Big\},
\]
where $\KL(q \| p) := \sum_{s'} q(s') \log(q(s')/p(s'))$ for probability vectors ($+\infty$ if
$q$ is not absolutely continuous with respect to $p$), and the threshold is $+\infty$ when
$n_h^t(s, a) = 0$ (the paper's $\alpha(0)/0 = \infty$): the constraint only concerns visited pairs.
\end{definition}

\textbf{Lean remark.} \texttt{eventKL X Y M δ : Set Ω} is
\texttt{\{ω | ∀ t h s a, klDiv (weightedMeasure (empTransAt X Y h s a t ω)) (M.trans h (s, a)) ≤ klThreshold S A H δ (visitCountAt X Y h s a t ω)\}}
with \texttt{klThreshold S A H δ n : ℝ≥0∞ := if n = 0 then ⊤ else ENNReal.ofReal (alphaKL S A H δ n / n)}
and Mathlib's \texttt{InformationTheory.klDiv} (with values in \texttt{ℝ≥0∞}, infinite when the
first measure is not absolutely continuous with respect to the second; on $\evKL$ every
visited pair has $\phat_h^t(s, a) \ll p_h(s, a)$). The paper's Appendix~A writes the threshold
$\alpha(n_h^t(s, a), \delta)$ without the division by $n_h^t(s, a)$; the confidence sets of its
Appendix~B and Lemma~19 have the division, which is the intended form.

\begin{definition}[Bernstein concentration event]\label{def:event_bern}
\lean{Essakine2026Tight.eventBern}\leanok
\uses{def:emp_trans, def:counts, def:rates, def:episodic_mdp}
Let $f = (f_h)_{h \le H + 1}$ with $f_h : \Ss \to \R$ and let $b = (b_h)_{h \le H}$ with $b_h \ge 0$
(in the applications, $b_h$ bounds the oscillation $\max f_{h+1} - \min f_{h+1}$ of $f_{h+1}$).
\[
  \evBern(f, b) := \Big\{ \forall t, h, (s, a) \text{ with } n := n_h^t(s, a) > 0 :\
    \big| (\phat_h^t - p_h) f_{h+1}(s, a) \big|
    \le \sqrt{2 \Var_{p_h}(f_{h+1})(s, a)\, \frac{\alpha^*(n, \delta)}{n}} + 3 b_h \frac{\alpha^*(n, \delta)}{n} \Big\} .
\]
\end{definition}

\begin{remark}[Deviation 12: step-dependent ranges, non-strict inequality]\label{rem:empirical_bern_event}
The paper's event $\mathcal E_f$ uses a single range $b$ for all steps and a strict inequality.
Two changes. (a) The ranges are step-dependent: Lemmas~7 and~12 of the paper
(Lemmas~\ref{lem:concentration_pos} and~\ref{lem:concentration_neg}) apply the event to
$f = Z^*$ with the range $e^{\beta(H - h)}$ (resp.\ $1 - e^{\beta(H - h)}$) of $Z^*_{h+1}$ \emph{at
the step $h$}, which is what makes the bonus of step $h$ scale with $e^{\beta(H - h)}$; a
uniform $b = e^{\beta H}$ would lose this. (b) The inequality is non-strict: with a strict
inequality the event would be \emph{empty} in the case $\beta < 0$, because at the step
$h = H$ the function $Z^*_{H+1} \equiv 1$ is constant and $b_H = 1 - e^{0} = 0$, so both sides
vanish for every visited pair. The concentration inequality (Corollary~\ref{cor:bernstein_explicit})
gives the non-strict form, and the analysis only uses the event as an upper bound.
\end{remark}

\textbf{Lean remark.} \texttt{eventBern X Y M δ f b : Set Ω} with \texttt{f : Fin (H + 1) → S → ℝ}
and \texttt{b : Fin H → ℝ}: for all \texttt{t h s a} with \texttt{0 < visitCountAt X Y h s a t ω},
\texttt{|vecExp (empTransAt …) (f h.succ) - vecExp (M.transVec h s a) (f h.succ)| ≤ √(2 * vecVar (M.transVec h s a) (f h.succ) * (alphaStar … n / n)) + 3 * b h * (alphaStar … n / n)}.

\begin{definition}[Counts concentration event]\label{def:event_cnt}
\lean{Essakine2026Tight.eventCnt}\leanok
\uses{def:counts, def:pseudo_counts, def:rates}
\[
  \evCnt := \Big\{ \forall t \in \N,\ \forall h \le H,\ \forall (s, a) :\
    n_h^t(s, a) \ge \tfrac12\, \nbar_h^t(s, a) - \alpha^{\mathrm{cnt}}(\delta) \Big\} .
\]
\end{definition}

\textbf{Lean remark.} \texttt{eventCnt M s₁ X Y δ : Set Ω}, with the random pseudo-counts of
Definition~\ref{def:pseudo_counts} (the outline's \texttt{eventCnt X Y δ P} loses the
dependence on \texttt{P}).

\begin{definition}[The good event]\label{def:good_event}
\lean{Essakine2026Tight.goodEvent}\leanok
\uses{def:event_kl, def:event_bern, def:event_cnt, def:opt_value, lem:opt_exp_value_range}
Let $Z^* = (Z^*_h)_{h \le H + 1}$ be the optimal exponential values of $\mathcal M$
(Definition~\ref{def:opt_value}). The \emph{good event} is
\[
  \evGood^+ := \evKL \cap \evBern(Z^*, b^+) \cap \evCnt \quad (\beta > 0),
  \qquad
  \evGood^- := \evKL \cap \evBern(Z^*, b^-) \cap \evCnt \quad (\beta < 0),
\]
with $b^+_h := e^{\beta(H - h)}$ and $b^-_h := 1 - e^{\beta(H - h)}$ for $h \le H$. By
Lemma~\ref{lem:opt_exp_value_range}, $Z^*_{h+1}$ takes values in $[1, e^{\beta(H - h)}]$ for
$\beta > 0$ and in $[e^{\beta(H - h)}, 1]$ for $\beta < 0$: an interval of length at most $b^\pm_h$,
so that $b^\pm_h \ge \max Z^*_{h+1} - \min Z^*_{h+1}$ as required by Lemma~\ref{lem:event_bern_prob}.
\end{definition}

\textbf{Lean remark.} \texttt{goodEvent M s₁ X Y β δ := eventKL X Y M δ ∩ eventBern X Y M δ (optExpValue M β) (fun h ↦ if 0 < β then Real.exp (β * k) else 1 - Real.exp (β * k)) ∩ eventCnt M s₁ X Y δ}
with \texttt{k = H - 1 - h} for the Lean step \texttt{h : Fin H} (the paper's $H - h$ for its
$1$-indexed step); both signs in one definition, the analysis chapters instantiating
\texttt{0 < β} or \texttt{β < 0}.

\section{Probability of the good event}
\label{sec:empirical_good_event}

The three probability bounds share one reduction: for a fixed triple $(h, s, a)$, the next
states observed at the successive visits of $(s, a)$ at step $h$ behave like an i.i.d.\
sample from $p_h(\cdot \mid s, a)$ (Lemma~\ref{lem:visits_iid}), and the empirical transition
after $t$ episodes is the empirical distribution of the first $n_h^t(s, a)$ of them
(Lemma~\ref{lem:emp_trans_reindex}). The time-uniform inequalities are then applied to this
sample, indexed by the number $n$ of visits, and transferred to all episodes $t$ at once.
Throughout this section, for a fixed $(h, s, a)$ we write $p := p_h(\cdot \mid s, a)$,
\[
  T_k := \min\{t \ge 1 : n_h^t(s, a) \ge k\} \in \N \cup \{\infty\}
  \quad\text{(the episode of the $k$-th visit)},
  \qquad
  W_k := s^{T_k}_{h+1} \text{ on } \{T_k < \infty\}
\]
(the next state observed at the $k$-th visit), and $\hat q_n(x_1, \dots, x_n) := \frac1n \sum_{k \le n} \delta_{x_k}$
for the empirical distribution of a finite sequence in $\Ss$. Each $T_k$ is a stopping time of
$(\mathcal F_t)$ (since $n_h^t(s, a)$ is $\mathcal F_t$-measurable), $T_k \le T_{k+1}$, and
$T_{n_h^t(s, a)} \le t$.

\begin{lemma}[Domination of the observed transitions by an i.i.d.\ sample, time-uniform form]\label{lem:empirical_visits_domination}
\uses{lem:visits_iid, def:counts, def:episode_env}
Fix $(h, s, a)$ and let $\xi_1, \xi_2, \dots$ be i.i.d.\ with law $p = p_h(\cdot \mid s, a)$ on some
auxiliary probability space $(\Omega', \Prob')$. Then:
\begin{enumerate}
  \item[(i)] for every $n \ge 1$ and every set $B \subseteq \Ss^n$,
        $\Prob\big(T_n < \infty,\ (W_1, \dots, W_n) \in B\big) \le \Prob'\big((\xi_1, \dots, \xi_n) \in B\big) = p^{\otimes n}(B)$;
  \item[(ii)] for every sequence of sets $B_n \subseteq \Ss^n$, $n \ge 1$,
        \[
          \Prob\big(\exists n \ge 1 :\ T_n < \infty,\ (W_1, \dots, W_n) \in B_n\big)
          \le \Prob'\big(\exists n \ge 1 :\ (\xi_1, \dots, \xi_n) \in B_n\big) .
        \]
\end{enumerate}
\end{lemma}

\begin{proof}
\uses{lem:visits_iid, def:counts}
(i) Lemma~\ref{lem:visits_iid}(i) gives the inequality for product sets
$B = B^1 \times \dots \times B^n$: $\Prob(W_1 \in B^1, \dots, W_n \in B^n, T_n < \infty) \le \prod_i p(B^i)$.
Since $\Ss^n$ is finite, an arbitrary $B$ is the disjoint union of its singletons, which are
product sets, and both sides are additive in $B$.

(ii) First-entrance decomposition. For $n \ge 1$ let
$B'_n := \{x \in \Ss^n : x \in B_n \text{ and } (x_1, \dots, x_k) \notin B_k \text{ for all } k < n\}$
and $E_n := \{T_n < \infty, (W_1, \dots, W_n) \in B'_n\}$. On $\{T_n < \infty\}$ all $T_k$, $k \le n$,
are finite, so $E_n$ is the event that $n$ is the \emph{first} index at which
``$T_n < \infty$ and $(W_1, \dots, W_n) \in B_n$'' holds (if it holds at $k < n$ then $T_k < \infty$
and $(W_1, \dots, W_k) \in B_k$, excluded by $B'_n$; conversely the first index at which it holds
is in some $E_n$). Hence the $E_n$ are pairwise disjoint and their union is the event of the
left-hand side, so its probability is $\sum_n \Prob(E_n) \le \sum_n p^{\otimes n}(B'_n)$ by (i).
Under $\Prob'$, the events $\{(\xi_1, \dots, \xi_n) \in B'_n\}$ are, by the same argument, the
pairwise disjoint events ``$n$ is the first index with $(\xi_1, \dots, \xi_n) \in B_n$'', whose
union is $\{\exists n, (\xi_1, \dots, \xi_n) \in B_n\}$; so
$\sum_n p^{\otimes n}(B'_n) = \Prob'(\exists n, (\xi_1, \dots, \xi_n) \in B_n)$.
\end{proof}

\textbf{Lean remark.} The auxiliary space can be taken to be $\Ss^{\N}$ with the product measure
\texttt{Measure.infinitePi (fun \_ ↦ weightedMeasure p)} (Mathlib \texttt{Measure.infinitePi}),
and the events of (ii) are measurable as countable unions. The lemma is the bridge between
the i.i.d.\ statements of Chapter~\ref{chap:pre_concentration} and the run; it is the
episodic analogue of LML's \texttt{Online/Bandit/RewardByCountMeasure.lean} (the ``arm'' is
$(h, s, a)$ and the ``reward'' the next state), to be contributed to LML with
Lemma~\ref{lem:visits_iid}.

\begin{lemma}[Time-uniform KL concentration of the empirical distribution]\label{lem:empirical_kl_time_uniform}
\uses{lem:ville, lem:num_types}
Let $p$ be a probability vector on a finite set $\mathcal X$ with $m := \abs{\mathcal X} \ge 1$
elements, $\xi_1, \xi_2, \dots$ i.i.d.\ with law $p$, and $\hat q_n := \hat q_n(\xi_1, \dots, \xi_n)$.
For every $\delta' \in (0, 1)$,
\[
  \Prob\Big( \exists n \ge 1 :\ n\, \KL(\hat q_n \| p) \ge \log\frac1{\delta'} + (m - 1)\log(n + 1) + 1 + \tfrac12 \log n \Big) \le \delta' ,
\]
and in particular
$\Prob\big( \exists n \ge 1 : \KL(\hat q_n \| p) \ge (\log(1/\delta') + m \log(e(n + 1)))/n \big) \le \delta'$.
\end{lemma}

\begin{proof}
\uses{lem:ville, lem:num_types}
\emph{Support.} Almost surely all $\xi_k$ lie in the support $\mathcal X_0 := \{x : p(x) > 0\}$;
on this event $\hat q_n$ is supported in $\mathcal X_0$ and $\KL(\hat q_n \| p) = \sum_{x \in \mathcal X_0} \hat q_n(x)\log(\hat q_n(x)/p(x)) < \infty$.
Replacing $\mathcal X$ by $\mathcal X_0$ ($m_0 := \abs{\mathcal X_0} \le m$ elements) only decreases
the threshold, so we assume $p(x) > 0$ for all $x$. Write $N_n(x) := n \hat q_n(x) = \#\{k \le n : \xi_k = x\}$,
so that, with $0 \log 0 := 0$,
\[
  n\, \KL(\hat q_n \| p) = \sum_x N_n(x) \log\frac{N_n(x)}{n} - \sum_{k \le n} \log p(\xi_k)
\]
(since $\sum_{k \le n} \log p(\xi_k) = \sum_x N_n(x) \log p(x)$).

\emph{A mixture martingale.} Let $\nu$ be the uniform probability measure on the simplex
$\Delta := \{q \in [0, 1]^{\mathcal X} : \sum_x q(x) = 1\}$ (normalized Lebesgue measure on the
$(m - 1)$-dimensional simplex, whose Lebesgue measure is $1/(m - 1)!$; the point mass if $m = 1$)
and
\[
  M_n := \int_\Delta \prod_{k \le n} \frac{q(\xi_k)}{p(\xi_k)}\, \nu(dq)
  = \Big(\prod_{k \le n} p(\xi_k)\Big)^{-1} \int_\Delta \prod_x q(x)^{N_n(x)}\, \nu(dq),
  \qquad M_0 = 1 .
\]
For a fixed $q$, $L_n(q) := \prod_{k \le n} q(\xi_k)/p(\xi_k)$ is a nonnegative martingale for the
natural filtration $\mathcal G_n := \sigma(\xi_1, \dots, \xi_n)$ with $L_0(q) = 1$:
$\E[L_{n+1}(q) \mid \mathcal G_n] = L_n(q)\, \E[q(\xi)/p(\xi)] = L_n(q) \sum_x q(x) = L_n(q)$. By
Fubini for nonnegative integrands, $M_n = \int L_n(q)\, \nu(dq)$ is a nonnegative martingale with
$M_0 = 1$, and Ville's inequality (Lemma~\ref{lem:ville}) gives $\Prob(\exists n, M_n \ge 1/\delta') \le \delta'$.

\emph{Lower bound on $M_n$.} The Dirichlet integral gives, for nonnegative integers $N(x)$ with
$\sum_x N(x) = n$, $\int_\Delta \prod_x q(x)^{N(x)}\, \nu(dq) = (m - 1)!\, \prod_x N(x)! / (n + m - 1)!$
(the integral of $\prod_x q(x)^{a_x - 1}$ over the simplex with respect to Lebesgue measure is
$\prod_x \Gamma(a_x)/\Gamma(\sum_x a_x)$, by induction on $m$ with the Beta integral). Hence
\[
  \log M_n - n\, \KL(\hat q_n \| p)
  = \log\frac{(m - 1)!}{(n + m - 1)!} + \sum_x \big[\log N_n(x)! - N_n(x)\log N_n(x)\big] + n \log n .
\]
Now $N! \ge (N/e)^N$ gives $\sum_x [\log N_n(x)! - N_n(x)\log N_n(x)] \ge -\sum_x N_n(x) = -n$;
and $(n + m - 1)!/(m - 1)! = \binom{n + m - 1}{m - 1}\, n!$ with $\binom{n + m - 1}{m - 1} \le (n + 1)^{m - 1}$
(Lemma~\ref{lem:num_types}) and $n! \le e \sqrt n\,(n/e)^n$ (Stirling's upper bound). Therefore
\[
  \log M_n - n\, \KL(\hat q_n \| p)
  \ge -(m - 1)\log(n + 1) - 1 - \tfrac12 \log n - n\log n + n - n + n\log n
  = -(m - 1)\log(n + 1) - 1 - \tfrac12 \log n .
\]

\emph{Conclusion.} On the event $\{\forall n, M_n < 1/\delta'\}$, of probability at least
$1 - \delta'$, for every $n \ge 1$:
$n\, \KL(\hat q_n \| p) \le \log M_n + (m - 1)\log(n + 1) + 1 + \frac12 \log n < \log(1/\delta') + (m - 1)\log(n + 1) + 1 + \frac12\log n$.
The simplified form follows from $1 + \frac12 \log n \le \log(e(n + 1))$ and $\log(n + 1) \le \log(e(n + 1))$.
\end{proof}

\begin{remark}[Why a union bound over $n$ of Sanov's bound does not suffice]\label{rem:empirical_sanov_union}
The paper proves the KL event by applying its Lemma~19 (Sanov's bound for a fixed sample size,
Theorem~\ref{thm:sanov_hp}) with the confidence $\delta/(3SAH)$ and a union bound over
$(h, s, a)$; the union over the sample sizes $n$ (equivalently over the episodes $t$) is left
implicit. It cannot be done with the rate $\alpha(n, \delta) = \log(3SAH/\delta) + S \log(8e(n + 1))$:
Sanov's bound gives $\Prob(\KL(\hat q_n \| p) > \alpha(n)/n) \le (n + 1)^{S - 1} e^{-\alpha(n)} = \frac{\delta}{3SAH}\,(8e)^{-S}\,\frac{1}{n + 1}$,
whose sum over $n \ge 1$ diverges. A union bound over $n$ with the summable weights
$1/(n(n + 1))$ would require the larger rate $\log(3SAH/\delta) + (S + 1)\log(8e(n + 1))$, i.e.\
$S + 1$ in place of $S$ in every downstream constant. Lemma~\ref{lem:empirical_kl_time_uniform}
(the mixture-martingale argument behind the time-uniform KL bound used by
\cite{menard2021fast}, whose rate is exactly $\alpha$) is the correct route and keeps the paper's
rates; it uses Ville's inequality and the type count of Chapter~\ref{chap:pre_concentration}
in place of Theorem~\ref{thm:sanov_hp}.
\end{remark}

\textbf{Lean remark.} The statement is library material for \texttt{Essakine2026Tight/Mathlib/}
(then Chapter~\ref{chap:pre_concentration}): a sample \texttt{ξ : ℕ → Ω → 𝒳} with
\texttt{iIndepFun} and \texttt{HasLaw (ξ k) (weightedMeasure p)}, the empirical measure as a
\texttt{weightedMeasure} and \texttt{klDiv}. Ingredients: Mathlib's
\texttt{Real.pow\_div\_factorial\_le\_exp} for $N^N/N! \le e^N$, the Stirling bound
$n! \le e\sqrt n\,(n/e)^n$ from \texttt{Stirling.stirlingSeq'\_antitone} and
\texttt{Stirling.stirlingSeq\_one} ($n!/(\sqrt{2n}(n/e)^n)$ is antitone with value $e/\sqrt 2$ at
$1$), and the Dirichlet integral, which is not in Mathlib (prove it by induction on $m$ with
\texttt{Real.betaIntegral}, or define $\nu$ as the law of a normalized vector of i.i.d.\
exponential variables and compute the moments). The martingale is with respect to the
natural filtration of \texttt{ξ}, and Lemma~\ref{lem:ville} is used in the form
$\Prob(\sup_n M_n \ge 1/\delta') \le \delta'$.

\begin{lemma}[Probability of the KL event]\label{lem:event_kl_prob}
\uses{def:event_kl, lem:empirical_kl_time_uniform, lem:empirical_visits_domination, lem:emp_trans_reindex, lem:rates_props}
For $\delta \in (0, 1)$, $\Prob(\evKL) \ge 1 - \delta/3$.
\end{lemma}

\begin{proof}
\uses{def:event_kl, lem:empirical_kl_time_uniform, lem:empirical_visits_domination, lem:emp_trans_reindex, def:rates, def:counts, def:emp_trans}
Fix $(h, s, a)$, put $p := p_h(\cdot \mid s, a)$, $\delta' := \delta/(3SAH)$, and let
$B_n := \{x \in \Ss^n : \KL(\hat q_n(x) \| p) > \alpha(n, \delta)/n\}$ for $n \ge 1$.

\emph{Reduction to the visits.} Let $\omega \notin \evKL$: there are $t$, $h$, $(s, a)$ with
$\KL(\phat_h^t(s, a) \| p_h(s, a)) > \alpha(n, \delta)/n$ where $n := n_h^t(s, a)$; since the threshold is
$+\infty$ for $n = 0$, $n \ge 1$, hence $T_n \le t < \infty$ and, by Lemma~\ref{lem:emp_trans_reindex},
$\phat_h^t(\cdot \mid s, a) = \hat q_n(W_1, \dots, W_n)$: so $(W_1, \dots, W_n) \in B_n$. Therefore
\[
  \Omega \setminus \evKL \subseteq \bigcup_{(h, s, a)} \big\{\exists n \ge 1 : T_n < \infty,\ (W_1, \dots, W_n) \in B_n\big\} .
\]

\emph{The i.i.d.\ bound.} With $\xi_1, \xi_2, \dots$ i.i.d.\ of law $p$ on a finite set of $S$
elements, Lemma~\ref{lem:empirical_kl_time_uniform} with $m = S$ and the confidence $\delta'$
gives $\Prob'(\exists n \ge 1 : \KL(\hat q_n(\xi) \| p) \ge (\log(1/\delta') + S\log(e(n + 1)))/n) \le \delta'$,
and $\log(1/\delta') + S\log(e(n + 1)) = \log(3SAH/\delta) + S \log(e(n + 1)) \le \alpha(n, \delta)$
(Definition~\ref{def:rates}: $e(n + 1) \le 8e(n + 1)$). Hence
$\Prob'(\exists n \ge 1 : (\xi_1, \dots, \xi_n) \in B_n) \le \delta'$.

\emph{Transfer and union bound.} Lemma~\ref{lem:empirical_visits_domination}(ii) gives
$\Prob(\exists n \ge 1 : T_n < \infty, (W_1, \dots, W_n) \in B_n) \le \delta'$ for each $(h, s, a)$, and
the union bound over the $SAH$ triples gives
$\Prob(\Omega \setminus \evKL) \le SAH \cdot \delta/(3SAH) = \delta/3$.
\end{proof}

\begin{lemma}[Probability of the Bernstein event]\label{lem:event_bern_prob}
\uses{def:event_bern, cor:bernstein_explicit, lem:empirical_visits_domination, lem:emp_trans_reindex, lem:rates_props}
Let $\delta \in (0, 1)$, $f = (f_h)_{h \le H + 1}$ and $b = (b_h)_{h \le H}$ with
$b_h \ge \max_{s'} f_{h+1}(s') - \min_{s'} f_{h+1}(s')$ for every $h \le H$. Then
$\Prob(\evBern(f, b)) \ge 1 - \delta/3$.
\end{lemma}

\begin{proof}
\uses{def:event_bern, cor:bernstein_explicit, thm:bernstein_self_normalized, lem:empirical_visits_domination, lem:emp_trans_reindex, def:rates, def:emp_trans, def:counts}
Fix $(h, s, a)$ and put $p := p_h(\cdot \mid s, a)$, $g := f_{h+1}$, $\bar g := pg$,
$\sigma^2 := \Var_p(g) = p(g - \bar g)^2$, $\delta' := \delta/(3SAH)$, and
\[
  B_n := \Big\{x \in \Ss^n : \big|(\hat q_n(x) - p) g\big| > \sqrt{2 \sigma^2 \alpha^*(n, \delta)/n} + 3 b_h\, \alpha^*(n, \delta)/n\Big\} .
\]

\emph{Reduction to the visits.} Exactly as in Lemma~\ref{lem:event_kl_prob}: if
$\omega \notin \evBern(f, b)$, some $t, h, (s, a)$ with $n := n_h^t(s, a) \ge 1$ violate the inequality,
$T_n \le t < \infty$, and $\phat_h^t(\cdot \mid s, a) = \hat q_n(W_1, \dots, W_n)$
(Lemma~\ref{lem:emp_trans_reindex}), so $(W_1, \dots, W_n) \in B_n$; thus
$\Omega \setminus \evBern(f, b) \subseteq \bigcup_{(h, s, a)} \{\exists n \ge 1 : T_n < \infty, (W_1, \dots, W_n) \in B_n\}$.

\emph{The i.i.d.\ bound.} Let $\xi_1, \xi_2, \dots$ be i.i.d.\ with law $p$ and
$Y_k := g(\xi_k) - \bar g$. The $Y_k$ are i.i.d., centered, with $\E[Y_k^2] = \sigma^2$, and
$\abs{Y_k} \le \max g - \min g \le b_h$ since $\bar g \in [\min g, \max g]$. If $b_h = 0$ then $g$ is
constant, $(\hat q_n(x) - p) g = 0$ for every $x$, and $B_n = \emptyset$ (the right-hand side is
nonnegative). If $b_h > 0$, apply Corollary~\ref{cor:bernstein_explicit} (the explicit form of the
self-normalized Bernstein inequality, Theorem~\ref{thm:bernstein_self_normalized}) to the
sequence $(Y_k)$ with its natural filtration, the constant weights $w_k := 1 \in [0, 1]$ (trivially
predictable), the bound $b_h$ and the confidence $\delta'$: here
$S_n = \sum_{k \le n} Y_k = n\,(\hat q_n(\xi) - p) g$ and $V_n = \sum_{k \le n} \E[Y_k^2 \mid \mathcal G_{k-1}] = n\sigma^2$,
so with probability at least $1 - \delta'$, for all $n \ge 1$,
\[
  n\,\big|(\hat q_n(\xi) - p) g\big| \le \sqrt{2 n \sigma^2 L_n} + 3 b_h L_n,
  \qquad L_n := \log\frac{4e(2n + 1)}{\delta'} .
\]
Since $4e(2n + 1) \le 8e(n + 1)$, $L_n \le \log(8e(n + 1)) + \log(3SAH/\delta) = \alpha^*(n, \delta)$
(Definition~\ref{def:rates}); dividing by $n$ and using the monotonicity of the right-hand side
in $L_n$, on the same event $(\xi_1, \dots, \xi_n) \notin B_n$ for all $n \ge 1$. Hence
$\Prob'(\exists n \ge 1 : (\xi_1, \dots, \xi_n) \in B_n) \le \delta'$.

\emph{Transfer and union bound.} Lemma~\ref{lem:empirical_visits_domination}(ii) and the union
bound over the $SAH$ triples give $\Prob(\Omega \setminus \evBern(f, b)) \le SAH \cdot \delta/(3SAH) = \delta/3$.
\end{proof}

\begin{remark}[The index of the Bernstein inequality]\label{rem:empirical_bernstein_index}
The paper's proof of Lemma~5 applies the self-normalized Bernstein inequality to the process
indexed by the \emph{episodes} $t$, with the predictable weights $w_t = \indic\{(s^t_h, a^t_h) = (s, a)\}$;
this gives the deviation bound with the logarithmic term $\log(4e(2t + 1)/\delta')$ in the
episode index $t$, and the paper then ``uses $n_h^t(s, a) \le t$ and the monotonicity of the
logarithm'' to replace it by $\alpha^*(n_h^t(s, a), \delta)$, which goes in the wrong direction
($\log(4e(2t+1)) \ge \log(4e(2n+1))$). Indexing the process by the number of visits, as above
and as in \cite{menard2021fast}, gives the term in $n$ directly; the domination lemma
(Lemma~\ref{lem:empirical_visits_domination}) is what allows to use the i.i.d.\ version of the
inequality. The paper's proof also chooses $\delta_{h, s, a} = \delta/(SAH)$ in the union bound,
which gives $\delta$ rather than the claimed $\delta/3$; the rates of Definition~\ref{def:rates}
contain $\log(3SAH/\delta)$, i.e.\ the choice $\delta_{h, s, a} = \delta/(3SAH)$ used here.
\end{remark}

\begin{lemma}[Probability of the counts event]\label{lem:event_cnt_prob}
\uses{def:event_cnt, thm:bernoulli_concentration, lem:pseudo_counts_increments, lem:rates_props}
For $\delta \in (0, 1)$, $\Prob(\evCnt) \ge 1 - \delta/3$.
\end{lemma}

\begin{proof}
\uses{def:event_cnt, thm:bernoulli_concentration, lem:pseudo_counts_increments, def:pseudo_counts, def:counts, def:rates}
Fix $(h, s, a)$ and $\delta' := \delta/(3SAH)$. For $i \ge 1$ let $X'_i := \indic\{(s^i_h, a^i_h) = (s, a)\}$,
the visit indicator of the episode $i$, and consider the filtration $(\mathcal F_i)_{i \ge 0}$.
$X'_i$ is a function of $(X_{i-1}, Y_{i-1})$, which is part of $\hist_i$, hence $\mathcal F_i$-measurable;
and by Lemma~\ref{lem:pseudo_counts_increments}(iii) (applied with $t = i - 1$),
$P_i := \Prob(X'_i = 1 \mid \mathcal F_{i-1}) = p^{\pi^i}_h(s, a)$ almost surely, which is
$\mathcal F_{i-1}$-measurable (a function of $X_{i-1} = \pi^i$). Theorem~\ref{thm:bernoulli_concentration}
(the Bernoulli maximal inequality, Lemma~20 of the paper, from \cite{dann2017unifying}) applied
to $(X'_i)$, $(P_i)$ and $(\mathcal F_i)$ with the confidence $\delta'$ gives
\[
  \Prob\Big(\exists t \ge 1 : \sum_{i \le t} X'_i < \frac12 \sum_{i \le t} P_i - \log\frac1{\delta'}\Big) \le \delta' .
\]
Now $\sum_{i \le t} X'_i = n_h^t(s, a)$ (Definition~\ref{def:counts}), $\sum_{i \le t} P_i = \nbar_h^t(s, a)$
(Definition~\ref{def:pseudo_counts}) and $\log(1/\delta') = \log(3SAH/\delta) = \alpha^{\mathrm{cnt}}(\delta)$;
for $t = 0$ the inequality $n_h^0 = 0 \ge 0 - \alpha^{\mathrm{cnt}}(\delta)$ holds since
$\alpha^{\mathrm{cnt}}(\delta) \ge 0$ (Lemma~\ref{lem:rates_props}). So the constraint of $\evCnt$ for the
triple $(h, s, a)$ fails with probability at most $\delta'$, and the union bound over the $SAH$
triples gives $\Prob(\Omega \setminus \evCnt) \le \delta/3$.
\end{proof}

\textbf{Lean remark.} The filtration is only needed inside Theorem~\ref{thm:bernoulli_concentration};
its hypotheses are supplied as: \texttt{X' i} is measurable with respect to the $\sigma$-algebra
generated by \texttt{histAt X Y i}, and the conditional law of \texttt{X' i} given
\texttt{(histAt X Y (i - 1), X (i - 1))} is the Bernoulli law of parameter
\texttt{occupancy M (X (i - 1)) s₁ h s a} (Lemma~\ref{lem:pseudo_counts_increments}(iii) as a
\texttt{HasCondDistrib} statement, from Lemma~\ref{lem:episode_env_law}). The martingale
$\exp(\sum_{i \le t}(P_i/2 - X'_i))$ of the proof of Theorem~\ref{thm:bernoulli_concentration} is
adapted to $(\mathcal F_t)$.

\begin{lemma}[Lemma 5: the good event has probability at least $1 - \delta$]\label{lem:good_event}
\uses{def:good_event, lem:event_kl_prob, lem:event_bern_prob, lem:event_cnt_prob, lem:opt_exp_value_range}
For $\beta > 0$, $\Prob(\evGood^+) \ge 1 - \delta$, and for $\beta < 0$, $\Prob(\evGood^-) \ge 1 - \delta$.
\end{lemma}

\begin{proof}
\uses{def:good_event, lem:event_kl_prob, lem:event_bern_prob, lem:event_cnt_prob, lem:opt_exp_value_range}
By Lemma~\ref{lem:opt_exp_value_range}, for every $h \le H$ the function $Z^*_{h+1}$ takes values
in $[1, e^{\beta(H - h)}]$ if $\beta > 0$ and in $[e^{\beta(H - h)}, 1]$ if $\beta < 0$, so
$\max Z^*_{h+1} - \min Z^*_{h+1} \le \abs{e^{\beta(H - h)} - 1} \le b^\pm_h$ (for $\beta > 0$,
$e^{\beta(H-h)} - 1 \le e^{\beta(H - h)} = b^+_h$; for $\beta < 0$, $1 - e^{\beta(H-h)} = b^-_h$). Hence
Lemma~\ref{lem:event_bern_prob} applies to $(Z^*, b^\pm)$ and gives
$\Prob(\evBern(Z^*, b^\pm)) \ge 1 - \delta/3$; with Lemmas~\ref{lem:event_kl_prob}
and~\ref{lem:event_cnt_prob}, the union bound gives
$\Prob(\Omega \setminus \evGood^\pm) \le \delta/3 + \delta/3 + \delta/3 = \delta$.
\end{proof}

\textbf{Lean remark.} \texttt{measureReal\_goodEvent\_ge : 1 - δ ≤ P.real (goodEvent M s₁ X Y β δ)}
for any algorithm-environment sequence \texttt{(X, Y)} in \texttt{statesEnv M s₁}, with the
hypotheses $\delta \in (0, 1)$, \texttt{M.HasRewardFn r} and $r \in [0, 1]$ (needed for the ranges of
$Z^*$); it is applied in Chapters~\ref{chap:analysis_pos} and~\ref{chap:analysis_neg} to runs
of Entropic-BPI, but holds for every algorithm.
